
\section{Software Engineering Aspect}
\label{sec:software_engineering_aspect}

The proposed dCeNN--ELM--ASP framework is implemented as a modular Python pipeline whose stages mirror the methodological workflow. Data preprocessing, model training, threshold calibration, anomaly detection, symbolic reasoning, utility mapping, and reporting are separated through explicit artifact interfaces and Makefile-based orchestration \cite{repo_thesis_grid_anomaly}. This stage-oriented design supports reproducibility, inspection, and maintainability, in accordance with established practices in scientific and machine-learning software \cite{wilson2014best,wilson2017good,elmenreich:repoducibility, bousbiat:TII2022,sculley2015debt}.

\subsection{Pipeline Modularity as an Architectural Principle}

The repository is organized around a clear top-level structure comprising \texttt{configs/}, \texttt{data/}, \texttt{src/}, \texttt{rules/}, \texttt{artifacts/}, \texttt{reports/}, and \texttt{edge/} directories, with a top-level \texttt{Makefile} used for orchestration \cite{repo_thesis_grid_anomaly}. This structure reflects an explicit architectural distinction between:
\begin{itemize}
	\item \textbf{configuration} (\texttt{configs/}),
	\item \textbf{input data} (\texttt{data/}),
	\item \textbf{pipeline logic} (\texttt{src/}),
	\item \textbf{symbolic knowledge} (\texttt{rules/}),
	\item \textbf{learned state and calibrated objects} (\texttt{artifacts/}),
	\item \textbf{outputs and analyses} (\texttt{reports/}),
	\item \textbf{deployment bundle} (\texttt{edge/}).
\end{itemize}
This separation is more than cosmetic. It ensures that code, data, model state, and outputs remain disentangled, which simplifies both debugging and scientific auditability.

At the highest level, the implementation can be viewed as a directed workflow
\begin{equation}
	\mathcal{P} =
	\{s_0, s_1, \dots, s_n\},
	\label{eq:pipeline_set}
\end{equation}
where each stage $s_i$ transforms an input artifact set $\mathcal{A}_i$ into an output artifact set $\mathcal{A}_{i+1}$:
\begin{equation}
	s_i: \mathcal{A}_i \rightarrow \mathcal{A}_{i+1}.
	\label{eq:artifact_transform}
\end{equation}
This means that the pipeline is composed as a chain of deterministic artifact transformations rather than as a monolithic procedural script. The practical consequence is that each stage can be executed, inspected, or replaced independently so long as it preserves the agreed interface artifact.

\subsection{Stage-Oriented Script Design in \texttt{src/}}

The most visible implementation of modularity appears in the numbered Python scripts under \texttt{src/}. The numbering convention itself is significant. It imposes a readable execution logic and communicates intended pipeline order to the developer, reviewer, or future maintainer. In the current repository, the scripts are grouped conceptually into the following blocks \cite{repo_thesis_grid_anomaly}:
\begin{itemize}
	\item \textbf{00--02}: data engineering and preprocessing,
	\item \textbf{03--05}: neural training, threshold calibration, and statistical detection,
	\item \textbf{06--07}: finance mapping and ASP reasoning,
	\item \textbf{08--11}: evaluation, edge export, and event/timeline construction,
	\item \textbf{15--25}: benchmarking, visualization, veto analysis, decision-boundary inspection, utility comparison, and report-table generation.
\end{itemize}

The numbering convention makes the intended execution order visible and links each methodological stage to a concrete implementation unit. It therefore improves navigation and makes the correspondence between the thesis workflow and the software structure easier to inspect.

\subsection{Makefile Orchestration and Workflow Automation}

A particularly strong software-engineering choice in this project is the use of a top-level \texttt{Makefile} to orchestrate the pipeline \cite{repo_thesis_grid_anomaly}. Rather than relying on manual execution of loosely connected scripts, the workflow is encoded as an explicit sequence of reproducible stages, which is consistent with prior work highlighting automation and modular pipeline design as key enablers of reliable machine-learning experimentation \cite{bousbiat:TII2022}, the implementation defines named targets such as:
\begin{itemize}
	\item \texttt{make holidays},
	\item \texttt{make preprocess},
	\item \texttt{make split},
	\item \texttt{make train},
	\item \texttt{make thresholds},
	\item \texttt{make detect},
	\item \texttt{make asp},
	\item \texttt{make finance},
	\item \texttt{make eval},
	\item \texttt{make edge}.
\end{itemize}
The full workflow is then reproducibly executed through
\begin{equation}
	\texttt{make all}.
\end{equation}

This is an important implementation decision because orchestration logic is often neglected in academic prototypes. In the present system, the \texttt{Makefile} acts as an executable workflow specification. It documents stage order, makes the project easier to reproduce on another machine, and reduces the chance of accidental omission or incorrect sequencing. The file also captures a subtle but meaningful design refinement: ASP reasoning is executed before finance mapping, ensuring that utility analysis is based on refined anomalies rather than raw neural output. This is a good example of software architecture reflecting methodological intent rather than merely wrapping code for convenience.

\subsection{Artifact-Based Interfaces Between Stages}

One of the strongest aspects of the implementation is the use of persisted artifacts as explicit interfaces between stages. Instead of relying on transient objects hidden within a single runtime, each important stage writes stable artifacts to disk. Examples include:
\begin{itemize}
	\item engineered datasets in \texttt{data/},
	\item supervised split matrices and scalers in \texttt{artifacts/},
	\item trained encoder weights and ELM heads in \texttt{artifacts/},
	\item calibrated thresholds in serialized configuration artifacts,
	\item anomaly tables and refined anomaly tables in \texttt{reports/tables/},
	\item model bundles for deployment in \texttt{edge/}.
\end{itemize}

This artifact-based design improves the software system in several ways.

First, it supports \textbf{traceability}. The researcher can inspect any intermediate product, such as the anomaly table before and after ASP refinement.

Second, it supports \textbf{restartability}. If a late-stage script fails, the full pipeline does not need to be rerun from scratch provided the required artifacts already exist.

Third, it supports \textbf{comparability}. Different model or threshold settings can be compared through distinct artifacts without entangling state across runs.

Fourth, it supports \textbf{extensibility}. A new model or refinement step can be introduced by consuming the same upstream artifact and producing the same downstream interface.

In formal terms, if stage $s_i$ expects artifact schema $\Sigma_i$, then modularity requires that any replacement stage $s_i'$ preserve that schema:
\begin{equation}
	s_i'(\mathcal{A}_i) \in \Sigma_{i+1}.
	\label{eq:schema_contract}
\end{equation}
This is exactly the software-engineering logic that makes the pipeline replaceable without becoming fragile.

\subsection{Decoupling Learning, Reasoning, and Application Logic}

Another important architectural property is the decoupling of the major reasoning strata of the system. The pipeline separates:
\begin{enumerate}
	\item \textbf{learning logic}, implemented in the dCeNN--ELM training and detection scripts;
	\item \textbf{symbolic logic}, implemented in the ASP rules and solver application stage;
	\item \textbf{application logic}, implemented in the finance-mapping and reporting stages.
\end{enumerate}

The separation is particularly important in a neuro-symbolic system. Symbolic knowledge is externalized in the \texttt{rules/} directory and evaluated using \texttt{clingo}, rather than being embedded as Python conditions within the learning code. Finance-aware utility mapping is likewise kept separate from anomaly generation. The software layout therefore preserves distinct roles for learning expected behaviour, evaluating rule consistency, and measuring downstream usefulness.

\subsection{Configuration-Driven Rather Than Hard-Coded Design}

Important experiment settings, including feature definitions, model parameters, threshold policies, and file paths, are stored in YAML files under \texttt{configs/}. This separates experimental settings from executable logic, reduces hard-coded duplication, and allows the same scripts to be reused under different configurations. The configuration structure is discussed in detail in Section~\ref{sec:hyperparameter_optimization}.

\subsection{Reporting, Benchmarking, and Pipeline Completeness}

The modular workflow also includes dedicated stages for benchmarking, visualization, veto analysis, decision-boundary inspection, utility comparison, and report-table generation. The implementation therefore extends reproducibility beyond model training to evaluation, interpretation, and the creation of thesis-ready outputs.

\begin{table}[H]
	\centering
	\caption{Software-engineering interpretation of the main repository modules.}
	\label{tab:software_architecture_modules}
	\small
	\begin{tabular}{|p{3.5cm}|p{5.0cm}|p{6.0cm}|}
		\hline
		\textbf{Repository module} & \textbf{Primary responsibility} & \textbf{Software-engineering significance} \\
		\hline
		\path{configs/} & YAML-based parameter and path management & Separates configuration from executable logic \\
		\hline
		\path{data/} & Raw and engineered datasets & Isolates input-state management from algorithms \\
		\hline
		\path{src/} & Stage-wise Python pipeline scripts & Enforces separation of concerns and readable workflow structure \\
		\hline
		\path{rules/} & ASP knowledge base & Externalizes symbolic reasoning from Python ML code \\
		\hline
		\path{artifacts/} & Trained model state, scalers, thresholds & Provides stable inter-stage interfaces and restartability \\
		\hline
		\path{reports/} & Tables, figures, and analysis outputs & Supports reproducible reporting and thesis-ready outputs \\
		\hline
		\path{edge/} & Lightweight deployment bundle & Decouples inference deployment from training environment \\
		\hline
		\path{Makefile} & End-to-end orchestration & Encodes workflow order and reproducibility in executable form \\
		\hline
	\end{tabular}
\end{table}

\subsection{Strengths and Limits of the Current Software Architecture}

The main strengths of the architecture are its stage modularity, explicit artifact interfaces, reproducible orchestration, separation of neural and symbolic logic, and integration of evaluation and reporting.

At the same time, a balanced assessment should acknowledge its limits. The project is still a research pipeline rather than a production-grade distributed service. Error handling is largely stage-local rather than centrally orchestrated. Artifact versioning is file-based rather than managed by a dedicated experiment-tracking system. The stage numbering is highly readable, but it also reflects a script-driven workflow rather than a formal workflow engine such as Airflow, Luigi, or Prefect. None of these observations weaken the thesis contribution; they simply clarify the level of engineering maturity being claimed. The implementation is best described as a \emph{well-structured scientific workflow system}, not as a large-scale enterprise software platform.

\subsection{Methodological Significance}

The software architecture supports the scientific reliability of the framework by mapping each methodological stage to a separate but connected implementation unit. Its modular scripts, persisted artifacts, externalized rules, and automated orchestration make the experimental workflow easier to reproduce, inspect, and extend.


\section{Edge Feasibility}
\label{sec:edge_feasibility}

Edge feasibility is evaluated as an explicit design objective of the proposed forecasting backbone. In this thesis, it has three dimensions: computational feasibility on modest CPU and memory resources, software feasibility through a minimal runtime stack, and deployment feasibility through a compact transferable model artifact. The dCeNN--ELM design supports these objectives through a compact latent representation, a linear analytical output head, and a NumPy-based runtime that does not require the complete PyTorch training environment \cite{repo_thesis_grid_anomaly}. This section examines these claims through hardware reporting and the edge-export mechanism implemented in \texttt{src/09\_edge\_export.py}.

\subsection{Hardware and Environment Transparency}

Edge-oriented performance claims must be interpreted relative to the hardware and software environment in which they were measured. The benchmark runs were executed in GitHub Codespaces with approximately \textbf{2 vCPU} and \textbf{7.8 GiB RAM}, using CPU inference to preserve comparison fairness. The corresponding \texttt{system\_hardware\_report.txt} records operating-system, processor, memory, Python, and library information, providing the computational context for the reported latency and model-size results.

The report does not present the target deployment environment as an ultra-microcontroller TinyML setting in the strictest sense. Instead, it positions the framework as suitable for \emph{lightweight CPU-first edge devices}, such as Raspberry Pi or Jetson-class platforms, where local anomaly detection is desirable but full cloud dependency is unnecessary or undesirable. This is consistent with the broader edge-AI literature, which emphasizes the importance of pushing intelligence closer to the data source when latency, resilience, or connectivity constraints matter \cite{antonini2023tinyml,howard2017mobilenets}.

\subsection{Why the dCeNN--ELM Core Is Naturally Lightweight}

The proposed forecasting backbone is especially well suited to edge deployment because of how its computational burden is distributed. Its inference-time logic consists of:
\begin{enumerate}
	\item feature standardization,
	\item one input projection into latent space,
	\item a small number of masked recurrent latent updates,
	\item one final linear output projection.
\end{enumerate}
If $\tilde{\mathbf{x}}_t \in \mathbb{R}^{d}$ denotes the standardized feature vector, the forward path can be written as
\begin{align}
	\mathbf{h}_t^{(0)} &= \tanh(W_{\mathrm{in}}\tilde{\mathbf{x}}_t + \mathbf{b}_{\mathrm{in}}), \\
	\mathbf{h}_t^{(k+1)} &=
	\tanh\!\left(
	(W_{\mathrm{cell}}\odot M)\mathbf{h}_t^{(k)} + \mathbf{b}_{\mathrm{cell}} + \gamma \mathbf{h}_t^{(k)}
	\right), \\
	\hat{\mathbf{y}}_t &= W_{\mathrm{out}}^{\top}\mathbf{h}_t^{(K)}.
	\label{eq:edge_forward}
\end{align}
The computational footprint remains limited because the latent dimension and number of recurrent steps are modest, while the block-structured mask restricts effective recurrent connectivity. The final forecast is produced by a linear output head, and the deployed forward pass uses NumPy rather than a complete training framework.

In asymptotic terms, if the input dimension is $d$, the latent dimension is $p$, the number of recurrent refinement steps is $K$, and the number of targets is $m$, then the dominant inference cost is approximately
\begin{equation}
	\mathcal{O}(dp + Kp^2 + pm).
	\label{eq:edge_complexity}
\end{equation}
Because $K$ and $p$ are small in the implemented configuration, the resulting cost remains modest and does not require backpropagation, attention over long sequences, or iterative optimization during inference.

\subsection{The Role of \texttt{09\_edge\_export.py}}

The practical mechanism that turns the trained forecasting model into a deployable edge artifact is the script \texttt{src/09\_edge\_export.py} \cite{repo_thesis_grid_anomaly}. This script is one of the most important implementation components for substantiating the edge-feasibility claim because it explicitly translates the training-time artifacts into a compact deployment bundle.

The script reads:
\begin{itemize}
	\item \texttt{configs/base.yaml},
	\item \texttt{configs/dcenn.yaml},
	\item the supervised training NPZ referenced by \texttt{npz\_path},
	\item \texttt{artifacts/scaler.npz},
	\item \texttt{artifacts/dcenn\_encoder.pt},
	\item \texttt{artifacts/elm\_heads.npz}.
\end{itemize}
It then writes two deployment outputs:
\begin{itemize}
	\item \texttt{edge/model\_bundle.npz},
	\item \texttt{edge/runtime\_infer.py}.
\end{itemize}
This separation is methodologically elegant. The first file is the compact parameter bundle; the second is the pure-NumPy runtime needed to execute inference on a target device.

\subsection{What Is Stored in the Edge Bundle}

The exported \texttt{model\_bundle.npz} is not merely a copy of training artifacts. It is a carefully selected subset of the trained state, stripped of unnecessary training-time objects and rewritten into edge-friendly arrays. Specifically, the bundle stores:
\begin{itemize}
	\item input projection weights and bias: \texttt{in\_W}, \texttt{in\_b};
	\item recurrent cell weights and bias: \texttt{cell\_W}, \texttt{cell\_b};
	\item the block-structured mask used in the dCeNN recurrence;
	\item model-shape metadata such as \texttt{steps}, \texttt{enc\_dim}, \texttt{block}, and \texttt{in\_dim};
	\item the ELM output matrix \texttt{W};
	\item feature-scaling statistics \texttt{scaler\_mean} and \texttt{scaler\_scale};
	\item \texttt{feature\_names} and \texttt{target\_names}.
\end{itemize}
This means the deployment bundle is self-contained with respect to forecasting inference: the edge device does not need the original training dataset, the PyTorch model class, or the scikit-learn \texttt{StandardScaler} object.

\begin{table}[H]
	\centering
	\caption{Main contents of the exported edge deployment bundle.}
	\label{tab:edge_bundle_contents}
	\small
	\begin{tabular}{|p{4.0cm}|p{4.5cm}|p{6.0cm}|}
		\hline
		\textbf{Bundle element} & \textbf{Stored object} & \textbf{Runtime purpose} \\
		\hline
		Encoder input weights & \path{in_W}, \path{in_b} & Maps standardized engineered features into latent space \\
		\hline
		Recurrent cell parameters & \path{cell_W}, \path{cell_b}, \path{mask} & Executes lightweight masked latent updates \\
		\hline
		Model metadata & \path{steps}, \path{enc_dim}, \path{block}, \path{in_dim} & Preserves architectural dimensions for runtime consistency \\
		\hline
		ELM output head & \path{W} & Produces multivariate forecasts from latent embeddings \\
		\hline
		Scaler statistics & \path{scaler_mean}, \path{scaler_scale} & Reproduces training-consistent standardization without scikit-learn \\
		\hline
		Feature/target metadata & \path{feature_names}, \path{target_names} & Supports dictionary-based inference and output interpretation \\
		\hline
	\end{tabular}
\end{table}

\subsection{Pure NumPy Runtime Inference}

A particularly strong edge-engineering choice is that \texttt{09\_edge\_export.py} automatically writes a companion runtime file, \texttt{edge/runtime\_infer.py}, containing a minimal NumPy-only implementation of the dCeNN--ELM forward pass \cite{repo_thesis_grid_anomaly}. This is crucial because it removes the need for PyTorch, scikit-learn, or any training-oriented framework on the target device. The runtime performs four operations:
\begin{enumerate}
	\item load the compressed bundle;
	\item reconstruct the standardized feature vector from a feature dictionary;
	\item execute the dCeNN latent forward pass;
	\item compute the output forecast through the stored ELM matrix.
\end{enumerate}

This design decouples the training and deployment environments. PyTorch is used during encoder training, whereas edge inference is reduced to NumPy matrix operations and \texttt{tanh} evaluations, improving portability across lightweight Linux-based devices.

\subsection{Why the ELM Head Improves Edge Feasibility}

The ELM component contributes to edge feasibility in a more subtle but equally important way. Since the output head is learned analytically and deployed as a simple weight matrix, the runtime does not need to execute a large decoder network or recurrent output head. Once the latent state $\mathbf{h}_t$ is available, inference reduces to
\begin{equation}
	\hat{\mathbf{y}}_t = \mathbf{h}_t^{\top}W,
	\label{eq:edge_linear_head}
\end{equation}
which is computationally trivial compared with the forward logic of many deep sequence models.

The ELM head therefore transfers the main fitting effort to the offline training stage, leaving only a compact and deterministic linear transformation for online inference.

\subsection{CPU-First Feasibility and Edge-Oriented Benchmark Interpretation}

The benchmark discussion is also important for interpreting edge feasibility. The reported environment of roughly \textbf{2 vCPU} and \textbf{7.8 GiB RAM} is not itself an embedded microcontroller setting, but it is a deliberately CPU-constrained and GPU-free benchmarking environment. That matters because the thesis claim is not that the model runs on a sub-milliwatt MCU; rather, it claims that the design is \emph{CPU-first}, \emph{lightweight}, and \emph{deployable on modest edge hardware}.

This is a reasonable and academically defensible claim. In the current edge-AI literature, there is a wide continuum between heavy cloud inference and ultra-constrained TinyML microcontrollers \cite{antonini2023tinyml}. The proposed model sits in the practical middle: it is intentionally much lighter than many deep recurrent or Transformer systems, yet still expressive enough for multivariate forecasting under weather uncertainty. The export script and pure NumPy runtime make that middle-ground position concrete.

\subsection{Strengths and Limits of the Current Edge Design}

The edge design provides a NumPy-only runtime, a compact self-contained model bundle, CPU-oriented inference, and a clear separation between training and deployment dependencies. However, the current export focuses on the forecasting core rather than the complete anomaly-detection pipeline. Threshold calibration, ASP reasoning, and finance mapping are not included in the exported edge runtime. The edge-readiness claim should therefore be interpreted as applying primarily to small CPU-based devices such as Raspberry Pi or Jetson-class platforms, rather than to bare-metal microcontrollers. Within this scope, the implementation demonstrates a credible lightweight forecasting backbone for edge-oriented anomaly-detection systems.

\subsection{Methodological Significance}

Edge feasibility is supported through the combination of compact model design, transparent CPU-based benchmarking, and an explicit export mechanism. \texttt{09\_edge\_export.py} packages the trained forecasting core into a self-contained NumPy bundle and minimal runtime, making the deployment claim concrete and reproducible.


\section{Hyperparameter Configuration and Selection}
\label{sec:hyperparameter_optimization}

The framework manages hyperparameters through YAML configuration files rather than embedding them throughout the source code. These files define the temporal split, feature construction, dCeNN architecture, ELM regularization, and residual thresholding policy \cite{repo_thesis_grid_anomaly}. The approach does not use a separate automated hyperparameter-search engine; instead, settings are adjusted through controlled, validation-informed experiments under a consistent pipeline. This is referred to here as \emph{configuration-managed parameter selection}. 

\subsection{Why Configuration Management Matters in This Thesis}

The pipeline contains coupled parameter families across data splitting, feature construction, representation learning, ridge fitting, threshold calibration, and edge export. For example, lag selection affects encoder input dimensionality, while threshold settings influence anomaly density and subsequent symbolic refinement. External configuration makes these interactions visible and separates experimental policy from execution logic, complementing the modular architecture described in Section~\ref{sec:software_engineering_aspect}.

\subsection{Configuration Files as Hyperparameter Modules}

The repository separates the main parameter families into \texttt{base.yaml}, \texttt{features.yaml}, \texttt{dcenn.yaml}, \texttt{elm.yaml}, and \texttt{thresholds.yaml}. These files respectively control the global experiment frame, temporal feature construction, encoder structure, ridge regularization, and residual-threshold policy. Their detailed parameters are presented in the following subsections and summarized in Table~\ref{tab:yaml_hyperparams}.

\subsection{Hyperparameter Families in Formal Terms}

The overall hyperparameter set of the system can be written abstractly as
\begin{equation}
	\Theta =
	\Theta_{\mathrm{base}}
	\cup
	\Theta_{\mathrm{feat}}
	\cup
	\Theta_{\mathrm{dcenn}}
	\cup
	\Theta_{\mathrm{elm}}
	\cup
	\Theta_{\mathrm{thr}},
	\label{eq:theta_union}
\end{equation}
where each subset governs one logical stage of the pipeline. More concretely:
\begin{align}
	\Theta_{\mathrm{base}} &= \{\text{split bounds},\ \text{horizon},\ \text{artifact paths}\}, \\
	\Theta_{\mathrm{feat}} &= \{\text{lags},\ \text{rolling windows},\ \text{winsorization columns}\}, \\
	\Theta_{\mathrm{dcenn}} &= \{\text{enc\_dim},\ \text{steps},\ \text{block},\ \text{epochs},\ \eta,\ \lambda_w,\ \text{seed}\}, \\
	\Theta_{\mathrm{elm}} &= \{\lambda_{\mathrm{ridge}}\}, \\
	\Theta_{\mathrm{thr}} &= \{\text{method},\ \alpha,\ \text{bucket\_by}\}.
	\label{eq:theta_families}
\end{align}
This formalization makes clear that the project does not have a single monolithic optimization problem. Instead, it has a staged hyperparameter design problem distributed across multiple subsystems.

\subsection{Base and Data-Split Hyperparameters}

The \texttt{base.yaml} file is the root configuration layer of the project \cite{repo_thesis_grid_anomaly}. It specifies:
\begin{itemize}
	\item the engineered dataset path,
	\item the holiday CSV path,
	\item the destination NPZ artifact path,
	\item the train/validation/test split boundaries,
	\item the forecast horizon.
\end{itemize}
These are not merely convenience parameters. They shape the entire downstream experiment. For example, the split dates define the temporal generalization regime of the thesis, while the horizon parameter controls whether the supervised framing is same-step or future-step predictive. Because these settings affect every later stage, centralizing them in one global file improves consistency and prevents silent divergence across scripts.

\subsection{Feature-Space Hyperparameters}

The \texttt{features.yaml} file contains the parameters that control the size and temporal memory of the supervised feature matrix \cite{repo_thesis_grid_anomaly}. In the current implementation, it defines:
\begin{itemize}
	\item lag offsets: \texttt{[1, 2, 3, 6, 12, 24, 48, 72]},
	\item rolling windows: \texttt{[3, 6, 12, 24, 48]},
	\item a set of columns subject to winsorization.
\end{itemize}

These settings are critical because they determine what temporal information is directly exposed to the forecasting backbone. If the lag set is too short, the model may miss daily or multi-day dependencies. If it is too long, the feature space becomes unnecessarily wide and may increase noise or computational burden. Similarly, rolling windows provide smoothed historical context but also change the statistical geometry of the input space.

\subsection{dCeNN Hyperparameters and Capacity Control}

The dCeNN encoder is the most structurally rich component of the neural pipeline, and its hyperparameters are appropriately externalized in \texttt{dcenn.yaml} \cite{repo_thesis_grid_anomaly}. The main parameters are:
\begin{itemize}
	\item \texttt{enc\_dim = 48}: latent embedding dimension,
	\item \texttt{steps = 2}: number of masked recurrent refinement steps,
	\item \texttt{block = 8}: block size for the structured connectivity mask,
	\item \texttt{ae\_epochs = 3}: autoencoder training epochs,
	\item \texttt{lr = 0.001}: learning rate,
	\item \texttt{weight\_decay = 1e-4}: regularization on trainable weights,
	\item \texttt{subsample\_train = 20000}: optional training subsampling,
	\item \texttt{seed = 42}: random seed.
\end{itemize}

These parameters control three distinct aspects of the model.

First, they control \textbf{representational capacity}. The latent dimension and number of refinement steps govern how much nonlinear structure the encoder can model.

Second, they control \textbf{structural inductive bias}. The block size influences the sparsity pattern of recurrent connectivity and thereby affects how localized or entangled the latent dynamics become.

Third, they control \textbf{training behaviour}. Epoch count, learning rate, weight decay, and random seed affect convergence, regularization, and reproducibility.

\subsection{ELM and Threshold Hyperparameters}

Compared with the dCeNN encoder, the ELM head is deliberately simple. Its main tunable parameter is the ridge regularization constant \texttt{l2} in \texttt{elm.yaml} \cite{repo_thesis_grid_anomaly}. Even though this is only one scalar, it plays an important role. If the value is too small, the analytical head may overfit the latent design matrix; if too large, it may underfit and wash out useful nonlinear structure encoded by the dCeNN. The advantage of keeping this hyperparameter isolated in its own YAML file is clarity: the output-layer regularization policy can be changed independently of the encoder design.

The thresholding stage is also hyperparameterized through its own configuration file. In \texttt{thresholds.yaml}, the current implementation sets:
\begin{itemize}
	\item \texttt{method = conformal},
	\item \texttt{alpha = 0.90},
	\item \texttt{bucket\_by = hour}.
\end{itemize}
These choices govern how validation residuals are converted into anomaly thresholds \cite{repo_thesis_grid_anomaly}. They are especially important in a weather-driven anomaly-detection system because threshold policy strongly affects the density, timing, and trustworthiness of anomaly candidates. Again, centralizing these values in a dedicated file allows threshold experimentation without rewriting detection code.

\subsection{Script-Level Consumption of YAML Configuration}

The modularity of the YAML system is reinforced by the fact that individual scripts explicitly load only the configuration layers they need. For example, the training script \texttt{03\_train\_dcenn\_elm.py} loads \texttt{base.yaml}, \texttt{dcenn.yaml}, and \texttt{elm.yaml} at startup \cite{repo_thesis_grid_anomaly}. The threshold-calibration script \texttt{04\_calibrate\_thresholds.py} loads \texttt{base.yaml}, \texttt{features.yaml}, \texttt{thresholds.yaml}, and \texttt{dcenn.yaml} \cite{repo_thesis_grid_anomaly}. The edge-export script \texttt{09\_edge\_export.py} again loads \texttt{base.yaml} and \texttt{dcenn.yaml} in order to reconstruct the minimal inference bundle correctly \cite{repo_thesis_grid_anomaly}.

This selective loading reduces coupling because each stage consumes only the configuration layers relevant to its responsibility.

\subsection{Optimization Strategy: Controlled Search Rather Than Exhaustive Automation}

It is important to state honestly what form of optimization is and is not implemented. The project does not include a separate hyperparameter-search engine performing large-scale automated sweeps or Bayesian optimization. Instead, optimization is achieved through:
\begin{itemize}
	\item explicit externalization of tunable parameters,
	\item deterministic reruns under revised YAML settings,
	\item validation-driven comparison of forecasting and anomaly outcomes,
	\item practical consideration of edge-feasibility constraints.
\end{itemize}
This can be described as \emph{controlled configuration-space search}. It is lighter than automated HPO but more transparent and more reproducible in a thesis context.

This choice is defensible for several reasons. First, the hyperparameter space of the implemented system is relatively structured and low dimensional. Second, the thesis prioritizes reproducibility and interpretability of settings over raw search volume. Third, the project’s edge-readiness goal means that some ``better'' hyperparameters in a purely statistical sense may still be undesirable if they lead to heavier models or brittle deployment behaviour. Thus, hyperparameter selection here is inherently multi-objective: it balances forecasting quality, anomaly usefulness, and efficiency.

\begin{table}[H]
	\centering
	\caption{Main YAML configuration files and their hyperparameter roles.}
	\label{tab:yaml_hyperparams}
	\small
	\begin{tabular}{|p{3.5cm}|p{5.5cm}|p{5.5cm}|}
		\hline
		\textbf{Configuration file} & \textbf{Representative contents} & \textbf{Experimental role} \\
		\hline
		\path{base.yaml} & dataset paths, holiday path, NPZ path, split boundaries, horizon & Defines the global experiment frame and temporal evaluation regime \\
		\hline
		\path{features.yaml} & lags, rolling windows, winsorization columns & Controls supervised feature-space construction and temporal memory \\
		\hline
		\path{dcenn.yaml} & \path{enc_dim}, \path{steps}, \path{block}, \path{ae_epochs}, \path{lr}, \path{weight_decay}, \path{seed} & Controls encoder capacity, structure, and training dynamics \\
		\hline
		\path{elm.yaml} & ridge parameter \path{l2} & Controls regularization of the analytical forecasting head \\
		\hline
		\path{thresholds.yaml} & \path{method}, \path{alpha}, \path{bucket_by} & Controls residual-to-anomaly calibration policy \\
		\hline
	\end{tabular}
\end{table}

\subsection{Strengths and Limits of the Current Hyperparameter Strategy}

The YAML-based strategy is reproducible, modular, and auditable because the reported settings are stored in named configuration files rather than hidden in source code. Its main limitation is that parameter selection remains researcher-guided; the project does not include automated large-scale search, experiment-tracking dashboards, or multi-run scheduling. For the scope of this thesis, the approach provides a practical balance between experimental control and implementation complexity.

\subsection{Methodological Significance}

The configuration system makes data splits, feature design, encoder structure, ridge regularization, and threshold policy explicit and reproducible. It therefore supports disciplined parameter selection without requiring an automated hyperparameter-optimization platform.

